\documentclass[12pt, a4paper]{article}

% Pacchetti base
\usepackage[utf8]{inputenc}
\usepackage[T1]{fontenc}
\usepackage[italian]{babel}

% Margini e layout
\usepackage{geometry}
\geometry{a4paper, margin=2.5cm}

% Indice cliccabile e formattazione
\usepackage{hyperref}
\usepackage{enumitem}
\usepackage{titlesec}

% Formattazione dei link
\hypersetup{
    colorlinks=true,
    linkcolor=black,
    filecolor=magenta,      
    urlcolor=blue,
}

\title{\textbf{Piano di Project Management}\\ \large Progetto: Shipping-on-the-Air}
\author{Team di Gestione}
\date{\today}

\begin{document}

\maketitle
\tableofcontents
\newpage

\section{Descrizione dell'Approccio Utilizzato}
Questo documento racconta il ``dietro le quinte'' organizzativo del progetto \textit{Shipping-on-the-Air}. L'obiettivo non è mostrare il codice, ma spiegare il \textit{come} e il \textit{perché} delle scelte gestionali che abbiamo preso per trasformare l'idea di una rete di consegne con i droni in una soluzione software reale e strutturata. Abbiamo seguito le classiche fasi del Project Management, adattandole però a un contesto tecnico e moderno.

Per rendere la documentazione viva e facilmente consultabile, abbiamo strutturato il lavoro simulando in tutto e per tutto l'infrastruttura di un'azienda tech reale. Abbiamo quindi centralizzato il progetto su un repository condiviso. Questa scelta ci ha portato diversi vantaggi pratici:
\begin{itemize}[label=\textbullet]
    \item \textbf{Collaborazione senza barriere:} avere uno spazio cloud ci ha permesso di lavorare in team anche a chilometri di distanza, senza mai perdere l'allineamento.
    \item \textbf{Controllo totale (Versioning):} ogni modifica ai documenti o ai requisiti è tracciata. Sappiamo sempre chi ha cambiato cosa e possiamo tornare indietro se facciamo un errore.
    \item \textbf{Navigazione ipertestuale:} i vari documenti del progetto non sono file isolati, ma sono collegati tra loro tramite percorsi relativi, rendendo l'esplorazione logica e veloce.
    \item \textbf{Stesura efficiente:} usare linguaggi di markup (come Markdown e LaTeX) ci ha permesso di concentrarci sui contenuti, lasciando che fosse il sistema a preoccuparsi di un'impaginazione pulita e professionale.
\end{itemize}

\subsection*{Il nostro Stack di Strumenti}
Per gestire la comunicazione, le tempistiche e la stesura del progetto, ci siamo appoggiati a un set di strumenti ben preciso:
\begin{itemize}[label=\textbullet]
    \item \textbf{Google meets:} Effettuiamo i meeting remoti con il cliente e i brainstorming interni, analizzando insieme bisogni e requisiti.
    \item \textbf{Trello:} la bacheca virtuale (stile Kanban) che ci aiuta a visualizzare lo stato dei task, assegnarli e monitorarne l'avanzamento.
    \item \textbf{Visual Studio Code:} Ide su cui scriviamo il codice, con estensioni per il versioning e la collaborazione in tempo reale.
    \item \textbf{Git e GitHub Desktop:} gli strumenti alla base del nostro repository, indispensabili per sincronizzare il lavoro di squadra senza sovrascriverci a vicenda.
    \item \textbf{LaTeX:} scelto scrivere la documentazione tecnica, ci ha permesso di produrre PDF professionali e ben formattati.
\end{itemize}

\subsection{Elenco dei documenti}
Per orientarsi facilmente all'interno del repository, ecco una mappa dei documenti chiave che abbiamo prodotto durante il ciclo di vita del progetto.

\subsubsection{Documentazione delle aziende}

\textbf{Documenti fornitore:}
\begin{itemize}[label=\textbullet]
\item \href{https://raw.githubusercontent.com/schwein69/Project-Management/master/docs/TeamSkills/team_skills.pdf}
{Startup Team Skills}
una tabella che mappa in modo onesto le nostre competenze tecniche (hard skills) e relazionali (soft skills), fondamentale per capire a chi assegnare i task più critici.

\item \href{https://raw.githubusercontent.com/schwein69/Project-Management/master/docs/ProjectClassification/project_classification.pdf}
{Project Classification}
La nostra valutazione della tipologia del progetto, utile per decidere il modello di Project Management più adatto.

\item \href{https://raw.githubusercontent.com/schwein69/Project-Management/master/docs/ProjectModificationRequest/request_template.pdf}
{Change Request Template}
il modulo ufficiale che facciamo firmare al cliente quando ci chiede di aggiungere o modificare funzioni in corso d'opera, è sempre meglio avere il nero su bianco.

\end{itemize}

\textbf{Documenti committente:}
\begin{itemize}[label=\textbullet]
    \item \textit{Analisi Operativa e Calendario:} i dati forniti dal cliente sui ritardi causati dall'attuale logistica e le mappe delle zone aeree in cui i droni dovranno operare.
\end{itemize}

\subsubsection{Documentazione del progetto}
% TODO: Aggiungere l'elenco dei documenti specifici (WBS, Piano dei Rischi, ecc.) una volta completati.

\section{Introduzione}
In questa sezione si delineano le informazioni fondamentali relative all'esecutore, \textbf{TechSolution}, e al committente del progetto \textit{Shipping-on-the-Air}. L'obiettivo è fornire una panoramica chiara del contesto aziendale, analizzando le dinamiche tra le parti e giustificando le scelte metodologiche adottate. Vengono inoltre presentate le competenze del team di sviluppo e le politiche interne che regolano il flusso di lavoro.

\subsection{Azienda esecutrice: TechSolution}
L'esecutore del progetto è \textbf{TechSolution}, un'azienda informatica specializzata nella progettazione e sviluppo di soluzioni software innovative. Pur essendo una realtà giovane, composta attualmente da tre figure professionali neo-laureate, si distingue per l'adozione di tecnologie all'avanguardia. Sebbene non disponga ancora di un know-how consolidato in specifici settori di mercato, la forte propensione all'innovazione ha permesso a TechSolution di ottenere visibilità mediatica tramite l'incubatore in cui opera. Questo posizionamento, unito a un'offerta economica competitiva, è stato il fattore determinante nella scelta da parte del committente. L'azienda è attualmente in fase di espansione e prevede l'inserimento a breve termine di nuove figure tecniche specializzate.

\subsubsection{Startup Team}
Il nucleo operativo di TechSolution è formato da tre membri, i quali ricoprono congiuntamente ruoli di Project Management e sviluppo. Le competenze individuali sono mappate nel documento \href{https://raw.githubusercontent.com/schwein69/Project-Management/master/docs/TeamSkills/team_skills.pdf}{\textit{Startup Team Skills}}, che valuta sia le \textit{pre-skills} (competenze relazionali, di comunicazione e di leadership) sia le \textit{pro-skills} (competenze tecniche e ingegneristiche). Questa duplice valutazione è ritenuta essenziale per un'assegnazione efficace dei task, specialmente all'interno di un team ristretto dove la flessibilità è fondamentale.

La valutazione delle competenze avviene tramite una scala numerica da 1 a 5, ottenuta mediando i risultati di autovalutazioni con le valutazioni incrociate tra colleghi. Per i futuri inserimenti, TechSolution prevede di affinare questo sistema mediante test pratici e aggiornamenti continui basati sulle performance sul campo.

\subsubsection{Politiche aziendali e Metodologia Scrumban}
Per mantenere il team agile e pronto agli imprevisti, in TechSolution abbiamo adottato un modello operativo \textbf{Scrumban}. All'atto pratico, questo significa che sfruttiamo una lavagna visiva (stile Kanban) per gestire il flusso di lavoro e limitare le attività simultanee, mantenendo però i ritmi, le iterazioni e i brevi incontri di allineamento tipici di Scrum. Questa metodologia ci garantisce un controllo rigoroso sulle scadenze, senza farci perdere la flessibilità necessaria per integrare i cambiamenti richiesti dal progetto. Tutte le politiche aziendali descritte in questa sezione sono state strutturate proprio per supportare questo approccio, ponendo una forte enfasi sul miglioramento continuo (\textit{Kaizen}) attraverso l'analisi critica a fine lavoro.

\subsubsection{Orari lavorativi e Flessibilità}
L'azienda segue un orario lavorativo in presenza dal lunedì al venerdì, dalle 09:00 alle 19:00, con una pausa pranzo prevista tra le 12:00 e le 13:30. A supporto dell'agilità intrinseca al team e per favorire il benessere lavorativo, è consentito lo svolgimento del lavoro in modalità da remoto (smart working) per una giornata alla settimana, mantenendo comunque un'elevata flessibilità per adattarsi alle esigenze operative e alle necessità del committente.

\subsubsection{Modalità di lavoro e Mob Programming}
Nel rispetto della policy aziendale settimanale di smart working, le attività procedono in modo integrato sia in sede che a distanza, garantendo in ogni caso la presenza per gli incontri chiave con gli stakeholder. Per affrontare sfide tecniche complesse e favorire il trasferimento di conoscenze, TechSolution incentiva, se necessario, il \textit{Mob Programming}, applicato sia durante la fase di sviluppo codice sia nella redazione della documentazione. Questa pratica stimola il brainstorming e migliora la qualità complessiva dell'output.

\subsubsection{Continuous Feedback e Knowledge Sharing}
Per mantenere alto il coinvolgimento e garantire la qualità del lavoro, TechSolution promuove attivamente una cultura basata sul \textit{Continuous Feedback}. In pratica, l'azienda valorizza e incentiva il tempo che i membri del team dedicano alla \textit{peer-review} incrociata (sia per lo sviluppo del codice che per la stesura della documentazione). Questo scambio strutturato di feedback costruttivi favorisce la collaborazione, accelera la crescita tecnica individuale e assicura che le \textit{best practices} vengano condivise e assimilate naturalmente da tutto il gruppo.

\subsubsection{Documentazione Continua}
In linea con l'approccio Agile, la documentazione non è relegata alla fine del progetto, ma è un processo continuo e iterativo. TechSolution predilige l'auto-generazione della documentazione tecnica, archiviando il tutto in un repository condiviso.

\subsubsection{Gestione dei Meeting}
TechSolution adotta un approccio metodologico rigoroso per la gestione degli incontri aziendali, avvalendosi di template standardizzati. Al fine di garantire la tracciabilità documentale a fini storici, ogni riunione richiede la redazione di un verbale contenente data, partecipanti e temi trattati. L'effettuazione di un incontro è subordinata alla preventiva condivisione di un'agenda e si conclude con la formalizzazione delle decisioni assunte.

\paragraph{Project Kick-Off Meeting}
Costituisce l'evento formale che sancisce il passaggio del progetto alla fase esecutiva. TechSolution prevede un'unica sessione strutturata per allineare tutte le parti coinvolte, garantendo la massima efficacia operativa. Al fine di condurre un incontro produttivo, l'agenda istituzionale si articola sulle seguenti aree fondamentali:
\begin{itemize}
    \item \textbf{Introduzioni:} Presentazione formale per offrire ai membri del team e agli stakeholder l'opportunità di conoscersi e di comprendere chiaramente i rispettivi ruoli all'interno del progetto.
    \item \textbf{Panoramica e Obiettivi:} Definizione chiara del "perché" e del "cosa" del progetto. In questa fase si delineano il perimetro d'azione (\textit{scope}), i \textit{deliverable} previsti e i risultati attesi, anche attraverso l'analisi condivisa del \textit{Project Definition Statement} (PDS).
    \item \textbf{Tempistiche e Milestone :} Revisione rigorosa della schedulazione, delle scadenze critiche e delle fasi in cui si articolerà il ciclo di vita del progetto.
    \item \textbf{Ruoli e Responsabilità :} Assegnazione formale degli incarichi per stabilire con esattezza chi eseguirà determinate attività e come i diversi membri del team (composto da tre risorse) collaboreranno tra loro.
    \item \textbf{Comunicazione e Strumenti:} Fissazione dei protocolli comunicativi aziendali, delle modalità di reportistica (es. \textit{Daily Meeting}) e degli strumenti di \textit{Project Management} adottati (come Trello, GitHub, Microsoft Teams, Google meets, etc...).
    \item \textbf{Rischi e Vincoli :} Identificazione dei potenziali colli di bottiglia (es. limiti tecnologici, normativi o di \textit{budget}) e pianificazione preventiva delle relative strategie di mitigazione.
    \item \textbf{Sessione Q\&A :} Apertura di uno spazio finale dedicato a domande e chiarimenti, volto a dissipare qualsiasi dubbio e garantire un'assoluta chiarezza prima dell'avvio effettivo dei lavori.
\end{itemize}

\paragraph{Project Scoping Meeting}
Incontri finalizzati alla concertazione con il committente e alla definizione del perimetro di progetto. Tali sessioni sono preposte alla produzione di deliverable analitici quali: \textit{Conditions of Satisfaction} (CoS), \textit{Requirements Breakdown Structure} (RBS) e \textit{Project Overview Statement} (POS). 
Qualora si rendano necessarie sessioni multiple, queste vengono distanziate di una settimana per consentire l'elaborazione dei dati, concludendosi con una \textit{review} critica sull'efficacia dell'incontro.

\paragraph{Daily Meeting}
Incontri di allineamento a cadenza quotidiana (fissati alle ore 09:30, con un limite temporale stringente di 30 minuti), concepiti per superare l'inefficienza dei tradizionali resoconti serali. 
Si discute brevemente lo stato di avanzamento dei task, le criticità emerse e le eventuali necessità di supporto. L'agenda è strutturata in tre punti chiave:
\begin{itemize}
    \item \textbf{Cosa ho fatto ieri:} Breve resoconto delle attività completate, con eventuale menzione di ostacoli superati.
    \item \textbf{Cosa farò oggi:} Presentazione dei task pianificati per la giornata, con indicazione delle priorità e delle risorse necessarie.
    \item \textbf{Blocchi o criticità:} Segnalazione immediata di problemi che potrebbero compromettere il rispetto della schedula, con eventuale richiesta di supporto o intervento da parte del team. Questi vengono registrati in un report da parte per tenere traccia delle criticità e delle azioni intraprese per risolverle.
\end{itemize}

\paragraph{Problem Resolution Meeting}
Incontri convocati \textit{on-demand} a seguito di criticità inibenti emerse durante i \textit{Daily Meeting}. L'agenda è focalizzata sulla circoscrizione del problema, l'assegnazione di un responsabile (\textit{owner}) e la definizione dei criteri per la risoluzione. Sebbene secondo le best practices di settore suggerisca il coinvolgimento esclusivo dei diretti interessati, l'attuale composizione a tre membri favorisce la partecipazione dell'intero gruppo per massimizzare l'efficacia del \textit{brainstorming}.

\paragraph{Project Review Meeting}
Organizzati in concomitanza con il rilascio di specifiche \textit{milestone}, questi incontri sono dedicati all'ispezione dei risultati e all'adattamento del piano di lavoro. La valutazione delle metriche di avanzamento e l'ideazione di azioni correttive richiedono la presenza obbligatoria del committente (o di un suo delegato formale), congiuntamente al team di progetto e agli stakeholder rilevanti.

\subsubsection{Gestione delle Comunicazioni}
La politica di comunicazione aziendale stabilisce che ogni approvazione legata ai requisiti di progetto o all'accettazione dei \textit{deliverable} debba essere sempre formalizzata per iscritto e archiviata in formato digitale, al fine di assicurarne la completa tracciabilità storica.
I Project Manager fungono da interfaccia primaria con il committente. Al termine di ogni confronto, il team redige un documento di sintesi (come, ad esempio, una \textit{Change Request}), sul quale il cliente è tenuto a esprimere un'esplicita approvazione o rifiuto. Questo meccanismo è stato introdotto per evitare i rischi operativi e contrattuali derivanti dall'accettazione tacita (il cosiddetto "silenzio-assenso").
Per quanto riguarda le comunicazioni interne, si producono documenti formali solo per arricchire la \textit{knowledge base} aziendale o per mantenere lo storico del progetto. I problemi e gli ostacoli quotidiani (\textit{impediments}) vengono invece discussi e risolti a voce durante i \textit{Daily Meeting}. 
Infine, per aumentare l'efficienza, i report scritti vengono redatti \textbf{esclusivamente} per i \textit{task} in ritardo o critici. Le attività che procedono regolarmente secondo i piani non richiedono alcuna inutile reportistica aggiuntiva.

\subsection{Committente}
A commissionare il progetto \textit{Shipping-on-the-Air} è NexusLog S.p.A., un operatore di spicco a livello nazionale nel settore dei corrieri espressi e delle spedizioni su larga scala, con un modello di servizio trasporto rivolto ai privati. Oltre a gestire i volumi di traffico tradizionali, NexusLog ha identificato la necessità strategica di integrare un nuovo servizio di consegne ultra-rapide dedicato ai pacchi leggeri e di espandere il proprio raggio d'azione entrando nel mercato del \textit{food delivery} con soluzioni proprietarie. Per raggiungere questo obiettivo, l'azienda ha scelto di utilizzare i droni e di esternare lo sviluppo del software necessario a un partner esterno. Bisogna però sottolineare che, nonostante la spinta verso l'innovazione, il committente vorrebbe comunque mantenere al sicuro, cioè senza rischiare di compromettere, la propria reputazione consolidata nel settore della logistica tradizionale. Pertanto, hanno deciso di fare outsourcing solo la parte tecnologica, lasciando a sé stessi la gestione operativa e commerciale del servizio, e con un budget limitato in termini di costi di sviluppo per capire se ha effettivamente senso investire in questa nuova linea di business.

Per inquadrare rapidamente le dimensioni e la complessità dell'azienda, è sufficiente osservare i numeri della sua attuale rete operativa:
\begin{itemize}
    \item Oltre 10.000 spedizioni giornaliere gestite tra Milano, Torino e Bologna;
    \item Un'infrastruttura logistica formata da 4 centri di smistamento e 30 filiali distribuite tra le città;
    \item Una squadra di 500 dipendenti.
\end{itemize}

A guidare questa transizione tecnologica per conto del cliente c'è l'Ing. Marco Rossi, nel duplice ruolo di \textit{Project Manager} e responsabile commerciale. 

\subsection{Relazione Esecutore-Committente}
Affrontare un progetto così innovativo e pieno di incognite — come coordinare in sicurezza il volo di droni in città per consegnare cibo e piccoli pacchi — è una grande sfida tecnologica. Per questo motivo, NexusLog ha deciso di evitare i classici fornitori informatici, cercando invece un partner più rapido e flessibile nel mondo delle \textit{startup}.

In questo contesto, anche se non avevamo mai collaborato prima, le nostre strade si sono incrociate. TechSolution si è subito distinta come la scelta ideale: a convincere il cliente sono state le nostre ottime competenze nello sviluppo software e la grande flessibilità garantita dal nostro metodo di lavoro (\textit{Scrumban}). Per un progetto così all'avanguardia, infatti, NexusLog aveva assoluto bisogno di un alleato capace di adattarsi velocemente ai cambiamenti e di gestire gli imprevisti senza bloccare i lavori.

\newpage

\section{Scoping}
In questa sezione vengono documentate tutte le decisioni strategiche prese durante la fase iniziale di \textit{scoping}. Dopo una breve panoramica sulle dinamiche del primo contatto, gli argomenti saranno strutturati in base ai \textit{deliverable} prodotti in questo gruppo di processo. Inoltre, verrà descritta l'organizzazione pratica dei \textit{meeting}, riportando per ciascuno l'ordine del giorno, i partecipanti coinvolti e una sintesi di quanto discusso.

\subsection{Primo contatto e analisi della richiesta}
Il primo contatto formale tra il committente e il nostro team è avvenuto tramite email aziendale. Nel suo messaggio introduttivo, NexusLog ha presentato brevemente la propria rete logistica e ha espresso l'intenzione di creare il sistema \textit{Shipping-on-the-Air} per gestire il nuovo servizio di consegne ultra-rapide tramite droni. Come spesso accade in questa fase introduttiva, la richiesta iniziale era ad alto livello, priva di specifiche tecniche o dettagli operativi.

Applicando le nostre metodologie operative, TechSolution ha risposto tempestivamente proponendo di organizzare un \textit{Project Scoping Meeting}. L'obiettivo era chiaro: analizzare a fondo la richiesta, chiarire le ambiguità e valutare concretamente la fattibilità tecnica dell'idea. 

Incrociando le rispettive agende, l'incontro è stato fissato per lunedì 01 maggio alle ore 10:00, presso la sala riunioni virtuale di TechSolution. A rappresentare l'azienda cliente si è presentato Mario Rossi, nel suo ruolo di \textit{Project Manager} e referente principale dell'iniziativa.

\subsection{Project Scoping Meeting}
\subsection{Resources Breakdown Structure (RBS)}
\subsubsection{Raccolta dei requisiti}
\subsubsection{Costruzione (I Sottosistemi)}
% Sottosistemi estratti dall'abstract per Shipping-on-the-Air
\paragraph{Gestione Ordini (Delivery \& Order)}
\paragraph{Logistica e Flotta (Drone Fleet)}
\paragraph{Navigazione (Routing \& GPS)}
\paragraph{Tracciamento Live (Tracking)}
\paragraph{Notifiche (Notification)}
\paragraph{Pagamenti (Payment)}

\subsection{Scelta del PMLC Model}
% TODO: Giustificare la scelta di un modello Iterativo o Adattivo per gestire l'incertezza tecnologica e i feedback.

\subsection{Budget e tempi}
\subsection{Project Overview Statement (POS)}
\subsection{Classificazione del progetto}
\subsection{Scelta del Core Team}

\section{Planning}
\subsection{Joint Project Planning Session}
\subsection{Work Breakdown Structure (WBS)}
\subsection{Prioritizzazione dei task}
\subsection{Stima delle risorse necessarie}
\subsection{Stima della durata dei task}
\subsection{Corrispettivo e stima dei costi}
\subsection{Project Network Diagram}
\subsubsection{Scope Bank}

\subsection{Analisi dei Rischi (Risk Management)}
% TODO: Focus sulla mitigazione proattiva di rischi esterni e tecnologici (es. perdita segnale GPS, variazioni meteo, aggiornamenti normativi).

\subsection{Project Definition Statement (PDS)}
\subsection{Piano di qualità}
\subsection{Scelta dei criteri di accettazione}
\subsection{Scelta del Developer Team e del Client Team}

\section{Launching / Executing}
\subsection{Scelta del Team e Bilanciamento}
\subsection{Regole operative per il team}
\subsection{Assegnamento Risorse e raffinamento della schedula}
\subsection{Assegnamento Responsabilità}
\subsection{Gestione delle comunicazioni}
\subsection{Documentazione e testing}
\subsection{Aggiornamento Project Definition Statement (PDS)}
\subsection{Work Packages}
\subsection{Gestione del processo di modifica dello Scope}

\section{Monitoring \& Controlling}
\subsection{Rispetto della schedula}
\subsection{Sistemi di reporting}
\subsection{Gestione dello Scope Bank}
\subsection{Issues Log}
\subsection{Problem Escalation Strategy}

\section{Chiusura del progetto}
\subsection{Procedura di accettazione}
\subsection{Installazione dei deliverable (Deploy flotta e software)}
\subsection{Documentazione del progetto}
\subsection{Post-implementation Audit}

\section{Bibliografia}

\end{document}
